\section{Conclusions}\label{sec:conclusions}

The sample covariance of equity returns mixes a small number of well-estimated
signal eigenvalues with several hundred poorly estimated noise eigenvalues.
Replacing the noise eigenvalues with their average (the trace-preserving CRE
cleaning prescribed by the Marchenko--Pastur
analysis~\cite{laloux1999,bun2017}) reduces the condition number by a factor
of roughly $300$ on S\&P~500 data, and its low-rank-plus-identity form admits
an exact direct solve via the Woodbury identity.

This paper makes three algorithmic contributions and one empirical contribution.

\emph{Direct Woodbury portfolio solver}. The primal-dual active-set loop requires
an inner solve of $T_{0,\mathcal{A}}^{\mathrm{RMT}} w = b$ at each outer step.
The naive approach assembles the $n_a \times n_a$ active-set system matrix and
solves by Cholesky at cost $O(n_a^2 k + n_a^3)$ per step. The Woodbury
identity reduces this to $O(n_a k^2 + k^3)$ by working with the $k \times k$
correction matrix $W$ instead, cheaper in flops precisely when $k < n_a$
(Proposition~\ref{prop:crossover}), which holds throughout the long-only
equity regime. The measured wall-clock advantage is $5\times$ warm and
$19\times$ cold vs KKT-Cholesky on a $50$-point synthetic frontier sweep, and
more than $20\times$ cold on a real S\&P~500 minimum-variance solve. Against
interior-point solvers the advantage is two orders of magnitude (Clarabel
called natively) and over three orders of magnitude (CVXPY out of the box).
The returned weights agree with interior-point solutions to solver accuracy
and attain marginally lower objective values. The same Woodbury kernel powers
an exact critical-line sweep that traces all $144$ frontier breakpoints for
the cost of a $50$-point grid (Section~\ref{ssec:cla}).

\emph{Randomised SVD preprocessing}. Computing the top-$k$ eigenpairs of
$\hat\Sigma$ via a dense eigendecomposition costs $O(n^2 T + n^3)$ and requires
$O(n^2)$ storage. A randomised truncated SVD applied directly to the return
matrix $X$ computes the same eigenpairs in $O(nTk)$ time and $O(nk)$ storage;
the measured portfolio impact of the approximation is below $3$\,bp in any
weight, an order of magnitude below the $k \pm 1$ threshold uncertainty
(Section~\ref{sec:ksensitivity}). At $n=3000$, randomised preprocessing is
$25\times$ faster than dense, and the full pipeline (rSVD + Woodbury solve)
is $4\times$ faster than the KKT-Cholesky solve alone, with no
preprocessing charged to the latter. For daily rolling windows an incremental
rank-two update of the tracked eigenpairs replaces the re-sketch at a
$23\times$ lower per-day cost, with median portfolio drift of $7$\,bp over
$100$ days (Section~\ref{sec:preprocessing}).

\emph{Scope of the structural trick}. Ridge terms, diagonal loadings,
externally specified factor models, and linear equality constraints preserve
the Woodbury structure (Remark~\ref{rem:ridge}); nonlinear shrinkage and
general rotationally invariant estimators~\cite{ledoitpeche2011,ledoit2012,bun2017}
do not, because they modify the full spectrum. CRE clipping is the member of
the RMT cleaning family that trades a known statistical refinement for an
exploitable algebraic structure.

\emph{Empirical validation}. The MP threshold is an insensitive
hyperparameter: misjudging $k$ by one factor moves portfolio weights by tens
of basis points and volatility by less than one basis point, and the
three-point audit costs milliseconds. Rolling out-of-sample backtests on two
universes, the high-dimensional S\&P~500 ($n/L \approx 0.98$) and the
well-sampled FTSE~100 ($n/L \approx 0.17$), show the CRE portfolio within
$0.1$ volatility points of linear shrinkage at the practitioner and oracle
intensities (and lowest outright on the FTSE), with all minimum-variance
strategies roughly a third less volatile than equal weight: the estimator
that enables the fast solver carries no measurable out-of-sample penalty in
either regime.

\paragraph{Practical recommendations}
For $n \leq 500$, either preprocessing path is practical; use dense if simplicity
matters. For $n > 1000$, use randomised SVD. For any sweep of $S \geq 2$ solves
(efficient frontier, rebalance sequence, sensitivity audit), the
preprocessing cost amortises to $O(nTk/S)$ per solve. The Woodbury solver is
preferred over KKT-Cholesky whenever $k < n_a$, which holds throughout the
parameter regime of long-only institutional equity portfolios.
